\documentclass{cswposter}
\usepackage{array}
\ExecuteBibliographyOptions{maxnames=2, minnames=1, giveninits=true}
\addbibresource{references.bib}
\newcommand{\img}[1]{\includegraphics[width=\linewidth]{#1}}

\begin{document}

\title{ActivePaws: A Wearable-Controlled Exergame Platform for Pediatric Motor Rehabilitation}
\author{\textbf{Marianna Pizzo}, Fabio Lazzari, Andrea De Filippis, Jacopo Roman\`o, Lorenzo Garavaglia, Simone Pittaccio, \textbf{Manuela Chessa}}

\begin{cswposter}

% --- COLUMN 1 ---

\posterbox[adjusted title = Motivation \& Research Questions]
{name = motivation, column=1}{
The \textit{Convention on the Rights of the Child} (UN, 1989) states that \textit{``Every child has a right to play''}. Yet children with upper limb neuromotor disorders often cannot access mainstream games, as conventional controllers exceed their motor abilities \cite{maggiorini2017evolution}. Research exergames address this but tend to be repetitive, poorly engaging, and graphically minimal compared to the games their healthy peers play \cite{biffi2017immersive}.

\textbf{ActivePaws} was designed to address both demands. On one hand, it uses a soft wearable orthosis as a controller~\cite{lazzari2026design} with \textbf{intuitive gestures} and \textbf{configurable alternative mappings} to accommodate different motor profiles. On the other hand, it provides a suite of exergames co-designed with therapists to be \textbf{visually and sonically engaging}, while remaining uncluttered enough to keep the motor task in focus.

{\color{color1}\textbf{RQ1.}} Is the system usable, acceptable, and engaging for children?\\
{\color{color1}\textbf{RQ2.}} Do alternative movement mappings affect game performance?\\
{\color{color1}\textbf{RQ3.}} Is there an optimal movement sensitivity threshold that facilitates successful gameplay?
}

\posterbox[adjusted title = The ActivePaws Exergame Suite]
{name = suite, column=1, below=motivation}{
ActivePaws is a wearable-controlled exergame suite for pediatric upper limb rehabilitation. A soft wrist orthosis, \textbf{Playcuff} \cite{lazzari2026design}, connects wirelessly via Bluetooth and serves as the game controller.

\medskip
The suite comprises three games. Each game offers \textbf{configurable alternative movement mappings} for key in-game actions, independently selectable before each session: one set is slightly more ecologically intuitive, ensuring accessibility across different residual motor abilities. Each game also provides a \textbf{calibration scene} where the therapist sets the sensitivity threshold for specific movements. A \textbf{two-phase tutorial}, combining illustrated slides with an interactive gameplay phase guided by voice-over, adapts to the configured movement set and accommodates children of different ages and reading abilities.

\medskip
{\color{color1}\textbf{Swift Rescuer:}} a vegan hunter frees trapped animals by shooting at them. A stability-to-shoot mechanic allows successful aim only when the crosshair is green. Four levels, the last introducing a cognitive sequencing challenge.

\smallskip
{\color{color1}\textbf{Sledding Adventures:}} the child steers a sled down a snowy track collecting items and avoiding obstacles, training forearm movements across four progressively demanding levels.

\smallskip
{\color{color1}\textbf{Galaxy Drift:}} the child pilots a spaceship against enemies while managing lives, fuel, and ammunition, training forearm and hand movements with a resource management cognitive layer.
}




\posterbox[adjusted title = Experimental Design]
{name = design, column=1 , below=suite}{
12 typically developing children (aged 4--12) participate in a single 120-minute session. Each child plays all three games twice, with the alternative movement configuration swapped between sessions. Assignment order is randomized to counterbalance exposure effects.

\smallskip
Data collected include in-game performance metrics and wristband movement classifications (automatically logged by the system), usability and engagement questionnaires (UAS~\cite{spinoso2026psychometric} and UES-SF~\cite{OBRIEN201828} adapted), and movement data via motion capture cameras and markers.
\medskip
\begin{tikzpicture}[x=0.155\linewidth, y=1cm, font=\footnotesize]
  \draw[->, line width=1pt, color=color2] (0,0) -- (6.1,0);

  \fill[color1] (0.5,0) circle (12pt);
  \node[above, yshift=14pt]  at (0.5,0) {\textbf{T0}};
  \node[below, yshift=-14pt] at (0.5,0) {5min};
  \node[below, yshift=-38pt] at (0.5,0) {Consent};

  \fill[color1] (1.5,0) circle (12pt);
  \node[above, yshift=14pt]  at (1.5,0) {\textbf{T1}};
  \node[below, yshift=-14pt] at (1.5,0) {45min};
  \node[below, yshift=-38pt, align=center] at (1.5,0) {Calib.\ + Game\\ with first set};
  \node[above, yshift=40pt, color1, align=center] at (1.5,0) {\textbf{Session 1}};

  \fill[color1] (3.0,0) circle (12pt);
  \node[above, yshift=14pt]  at (3.0,0) {\textbf{T2}};
  \node[below, yshift=-14pt] at (3.0,0) {15min};
  \node[below, yshift=-38pt, align=center] at (3.0,0) {Break +\\ Questionnaires\\ on Session 1};

  \fill[color1] (4.5,0) circle (12pt);
  \node[above, yshift=14pt]  at (4.5,0) {\textbf{T3}};
  \node[below, yshift=-14pt] at (4.5,0) {45min};
  \node[below, yshift=-38pt, align=center] at (4.5,0) {Calib.\ + Game\\ with alt.\ set};
  \node[above, yshift=40pt, color1, align=center] at (4.5,0) {\textbf{Session 2}};

  \fill[color1] (5.5,0) circle (12pt);
  \node[above, yshift=14pt]  at (5.5,0) {\textbf{T4}};
  \node[below, yshift=-14pt] at (5.5,0) {10min};
  \node[below, yshift=-38pt, align=center] at (5.5,0) {End +\\ Questionnaires\\ on Session 2};

  \node[right] at (6.1,0) {\textbf{120 min}};
\end{tikzpicture}
}
\posterbox[adjusted title = Future Works]
{name = future, column=1, below=design}{
The system will be evaluated in studies with children with neuromotor disorders within the \textbf{Fit4MedRob} investigation, providing evidence of therapeutic efficacy, motor learning, and sustained engagement.}




% --- COLUMN 2 ---

\posterbox[adjusted title = Game Screenshots]
{name = screenshots, column=2}{
\begin{tabular}{m{0.05\linewidth} m{0.3\linewidth} m{0.3\linewidth} m{0.3\linewidth}}
\centering & \centering\textbf{Swift Rescuer} & \centering\textbf{Sledding Adventures} & \centering\arraybackslash\textbf{Galaxy Drift} \\[0.3em]
\centering\rotatebox{90}{\footnotesize Initial UI} & \img{assets/Figure/interfacciainiziale/sr.png} & \img{assets/Figure/interfacciainiziale/sa.png} & \img{assets/Figure/interfacciainiziale/gd.png} \\[0.2em]

\centering\rotatebox{90}{\footnotesize  Alt. Movements\hspace{2pt}} & \img{assets/Figure/movimenti/sr.png} & \img{assets/Figure/movimenti/sa.png} & \img{assets/Figure/movimenti/gd.png} \\[0.2em]

\centering\rotatebox{90}{\footnotesize Calibration\hspace{2pt}} & \img{assets/Figure/calibrazione/sr.png} & \img{assets/Figure/calibrazione/sa.png} & \img{assets/Figure/calibrazione/gd.png} \\[0.2em]

\centering\rotatebox{90}{\footnotesize Text. Tutorial\hspace{2pt}} & \img{assets/Figure/tutorial/sr.png} & \img{assets/Figure/tutorial/sa.png} & \img{assets/Figure/tutorial/gd.png} \\[0.2em]

\centering\rotatebox{90}{\footnotesize Inter. Tutorial\hspace{2pt}} & \img{assets/Figure/gameplaytu/sr.png} & \img{assets/Figure/gameplaytu/sa.png} & \img{assets/Figure/gameplaytu/gd.png} \\[0.2em]

\centering\rotatebox{90}{\footnotesize Gameplay\hspace{2pt}} & \img{assets/Figure/gioco/sr.png} & \img{assets/Figure/gioco/sa.png} & \img{assets/Figure/gioco/gd.png} \\

\end{tabular}
}


\posterbox[]
{name = childphoto, column=2, below=screenshots}{
\centering
\includegraphics[width=0.7\linewidth]{assets/img.png}
}



\posterbox[adjusted title = References]
{name = refs, column=2, below=childphoto}{
\printbibliography[heading=none]
}

\end{cswposter}

\begin{tikzpicture}[remember picture, overlay]
  \node[anchor=south, yshift=1.5cm, font=\fontsize{30}{36}\selectfont]
    at (current page.south)
    {\texttt{marianna.pizzo@edu.unige.it} -- \textbf{PILab}};
\end{tikzpicture}

\end{document}